Calculando as cargas atuantes devido ao vento direto na perna:

\begin{align*}
  A_p & = b \cdot H_s                 \\
  A_p & = \num{0.07} \cdot 8          \\
  A_p & = \boxed{\SI{0.56}{\meter^2}}
\end{align*}

A pressão aerodinâmica é de $P_a = \SI{250}{\pascal}$.

Com base na norma NBR 8400 Tabela 8 \cite{nbr8400}, obteve-se o fator $C = 2$.

\begin{align*}
  F_{V_A} & = A_p \cdot P_a \cdot C              \\
  F_{V_A} & = \num{0.56} \cdot \num{250} \cdot 2 \\
  F_{V_A} & = \boxed{\SI{280}{\newton}}
\end{align*}

\begin{align*}
  q_{V_A} & = \frac{F_{V_A}}{H_s}                \\
  q_{V_A} & = \frac{\SI{280}{\newton}}{8}        \\
  q_{V_A} & = \boxed{\SI{35}{\newton\per\meter}}
\end{align*}

Calculando as cargas atuantes devido ao vento lateral no pórtico:

\begin{align*}
  K & = \frac{H_s}{V} \\
  K & = \frac{8}{8}   \\
  K & = \boxed{1}
\end{align*}

\begin{align*}
  A_y & = B_y = \frac{q_V \cdot {H_s}^2}{2 V}          \\
  A_y & = B_y = \frac{\num{35} \cdot {8}^2}{2 \cdot 8} \\
  A_y & = B_y = \boxed{\SI{140}{\newton}}
\end{align*}

\begin{align*}
  A_x & = \frac{11 k + 18}{8 \left( 2 k + 3 \right)} \cdot q_V \cdot H_s                \\
  A_x & = \frac{11 \cdot 1 + 18}{8 \left( 2 \cdot 1 + 3 \right)} \cdot \num{35} \cdot 8 \\
  A_x & = \boxed{\SI{203}{\newton}}
\end{align*}

\begin{align*}
  B_x & = \frac{5 k + 6}{8 \left( 2 k + 3 \right)} \cdot q_V \cdot H_s                \\
  B_x & = \frac{5 \cdot 1 + 6}{8 \left( 2 \cdot 1 + 3 \right)} \cdot \num{35} \cdot 8 \\
  B_x & = \boxed{\SI{77}{\newton}}
\end{align*}

\begin{align*}
  M_c & = \frac{3 \left(k + 2 \right)}{8 \left( 2 k + 3 \right)} \cdot q_V \cdot {H_s}^2          \\
  M_c & = \frac{3 \left(1 + 2 \right)}{8 \left( 2 \cdot 1 + 3 \right)} \cdot \num{35} \cdot {8}^2 \\
  M_c & = \boxed{\SI{504}{\newton\meter}}
\end{align*}

\begin{align*}
  M_d & = \frac{5 k + 6}{8 \left( 2 k + 3 \right)} \cdot q_V \cdot {H_s}^2                \\
  M_d & = \frac{5 \cdot 1 + 6}{8 \left( 2 \cdot 1 + 3 \right)} \cdot \num{35} \cdot {8}^2 \\
  M_d & = \boxed{\SI{616}{\newton\meter}}
\end{align*}

\begin{align*}
  M_{\text{max}} & = \frac{{\left( 11 k + 18 \right)}^2}{128 {\left(2 k + 3 \right)}^2} \cdot q_V \cdot {H_s}^2                \\
  M_{\text{max}} & = \frac{{\left( 11 \cdot 1 + 18 \right)}^2}{128 {\left(2 \cdot 1 + 3 \right)}^2} \cdot \num{35} \cdot {8}^2 \\
  M_{\text{max}} & = \boxed{\SI{602.83}{\newton\meter}}
\end{align*}

\begin{align*}
  y & = \frac{11 k + 18}{8 \left( 2 k + 3 \right)} \cdot H_s           \\
  y & = \frac{11 \cdot 1 + 18}{8 \left( 2 \cdot 1 + 3 \right)} \cdot 8 \\
  y & = \boxed{\SI{5.8}{\meter}}
\end{align*}

Calculando as cargas atuantes devido ao vento na perna oposta:

Com base na norma NBR 8400 Tabela 9 \cite{nbr8400}, obteve-se o fator $\phi = 1$.

\begin{align*}
  F_{V_B} & = A_p \cdot P_a \cdot C \cdot \phi           \\
  F_{V_B} & = \num{0.56} \cdot \num{250} \cdot 2 \cdot 1 \\
  F_{V_B} & = \boxed{\SI{280}{\newton}}
\end{align*}

\begin{align*}
  q_{V_B} & = \frac{F_{V_B}}{H_s}                \\
  q_{V_B} & = \frac{\SI{280}{\newton}}{8}        \\
  q_{V_B} & = \boxed{\SI{35}{\newton\per\meter}}
\end{align*}
